\section{Construction}
\label{sec:construction}

\paragraph{Primitives and interactions.}
We describe local algorithms on typed inputs as \emph{primitives} and
the protocol's public-message rules as \emph{interactions}.
Each interaction has a \emph{When} test and a
\emph{What} computation, with the invoking pid as its first argument.
The protocol specifies the interaction calls and their arguments.
The ledger argument $L_r$ is the snapshot frozen at the start of round $r$.

Whenever the ledger invokes a party's step algorithm, the party evaluates
\emph{When} for each interaction call whose first argument is its pid.
If none returns $\top$, its message is $\bot$.
Otherwise, it computes \emph{What} for every call whose test returns $\top$.
These contributions are conveyed in a single posting.

For the computation interactions below, \emph{What} describes the
meaningful contribution, such as a share.  The contribution is recoverable
from the party's secret key and the interaction's public arguments.
The actual message returned by $\ComputationStep$ is one copy of $\sk$;
readers recover the contributions by recomputing \emph{What} with this
key and the same public arguments.  Thus, saying that a member sends its
share is shorthand for this key disclosure and recovery.  Application
interactions specify the actual messages to include in their posting.

\subsection{Common reference string}
\label{sec:construction-crs}

The ledger runs the $\CRS$ algorithm of Figure~\ref{fig:construction-crs}
at initialization and includes its output in the initial $\Setup$ record.
The resulting parameters $U$ are used for public sampling throughout
the protocol.

\begin{figure}[tbp]
\centering
\fbox{%
\begin{minipage}{0.95\textwidth}
\small
\SubprotocolHeading{\CRS}
\GadgetHeading{Primitives}
\begin{description}
  \item[$\CRS(1^\lambda)$.]
    $\lambda$: the security parameter.
    $\Setup^H$: the universal sampler's setup algorithm, using the
    public random oracle $H$ (Section~\ref{sec:universal-samplers}).
    Sample $U\sample\Setup^H(1^\lambda)$.
    Return $\crs=(U)$, where $U$ is the public sampler parameter string.
\end{description}
\end{minipage}}
\caption{Common reference string generation.}
\label{fig:construction-crs}
\end{figure}

\subsection{Conditional release}
\label{sec:conditional-release}

We begin with the committee mechanism used in
Section~\ref{sec:overview-gadgets}: a value is shared and the shares are
encrypted to committee members.  When a public condition holds, the
members reveal their decryption keys, allowing public reconstruction.
We call this gadget \emph{encrypted sharing with conditional release}
($\ESCR$) and use it for both nonces and garbled input labels.

Figure~\ref{fig:conditional-release} defines the preparation and recovery
algorithms and the release interaction.  In our construction,
$\ESCR.\Prepare$ runs inside the distribution used by the universal
sampler.  Its output supplies the public ciphertexts for the interaction.

\begin{figure}[tbp]
\centering
\fbox{%
\begin{minipage}{0.95\textwidth}
\small
\SubprotocolHeading{\ESCR}
Use robust sharing with parameters $(n,t)$, secret space $\mathcal M$,
and share space $\mathcal S$ (Section~\ref{sec:secret-sharing}), and RNCE
with a fixed padded encoding of shares (Section~\ref{sec:nce}).
All vectors below have $n$ indexed entries.
\GadgetHeading{Primitives}
\begin{description}
  \item[$\ESCR.\Prepare(v,\boldsymbol{\pk})\to\boldsymbol c$.]
    Input: $v\in\mathcal M$ and a vector
    $\boldsymbol{\pk}=(\pk_1,\ldots,\pk_n)$ of RNCE public keys.
    Output: a vector $\boldsymbol c=(c_1,\ldots,c_n)$ of ciphertexts.
    Sample $(v_1,\ldots,v_n)\sample\Share(v)$.
    For each $i\in[n]$, independently sample
    $c_i\sample\Enc(\pk_i,v_i)$, using the fixed share encoding.
    Return $\boldsymbol c=(c_1,\ldots,c_n)$.

  \item[$\ESCR.\Recover(\boldsymbol c,\widehat{\boldsymbol{\sk}})\to v$.]
    Input: $\boldsymbol c=(c_1,\ldots,c_n)\in(\bits^*)^n$ and
    $\widehat{\boldsymbol{\sk}}=(\widehat{\sk}_1,\ldots,\widehat{\sk}_n)
      \in(\bits^*\cup\{\bot\})^n$.
    Output: an element of $\mathcal M\cup\{\bot\}$.
    For each $i\in[n]$, set $\widehat v_i=\bot$ if
    $\widehat{\sk}_i=\bot$.  Otherwise compute
    $\Dec(\widehat{\sk}_i,c_i)$ and decode the result as a share
    $\widehat v_i\in\mathcal S$, using $\bot$ on decryption or decoding
    failure.
    Return $\Reconstruct(\widehat v_1,\ldots,\widehat v_n)$.
\end{description}

\GadgetHeading{Interactions}
\begin{description}
  \item[$\ESCR.\Release(\pk,\sk,L_r;\boldsymbol{\pk},\boldsymbol c,B)$.]
    $\pk,\sk$: the invoking party's RNCE public and secret keys;
    its pid is $\pk$.
    $L_r$: the ledger snapshot frozen at round start.
    $\boldsymbol{\pk}$: an ordered committee of distinct registered
    computation pids, each reserved for this use.
    $\boldsymbol c$: the ciphertext vector for this committee.
    $B$: a deterministic polynomial-time predicate on ledger snapshots,
    supplied by the calling protocol.
    \begin{description}
      \item[When.]
        $\pk\in\boldsymbol{\pk}$, $\pk\in\Schedule(L_r)$,
        and $B(L_r)=1$.
      \item[What.]
        For the unique $i$ with $\pk=\pk_i$, decrypt $c_i$ using $\sk$
        and return the decoded share $v_i$, or $\bot$ on decryption
        or decoding failure.  The actual ledger posting is $\sk$.
    \end{description}
\end{description}

\end{minipage}}
\caption{$\ESCR$: encrypted sharing with conditional release.}
\label{fig:conditional-release}
\end{figure}

The protocol reserves the committee and schedules its members together
for a release round $r$ whose snapshot satisfies $B(L_r)=1$.
At the next snapshot $L_{r+1}$, each reader builds
$\widehat{\boldsymbol{\sk}}$ by taking $\widehat{\sk}_i$ from
$(\Message,r,\pk_i,\widehat{\sk}_i)$, with $\bot$ if no such record
occurs.  These values, together with the prepared ciphertext vector,
are the inputs to $\ESCR.\Recover$.

\subsection{Shared coins}
\label{sec:shared-coins}

Each nonce $Q_s$ is prepared by a separate universal sample and held by
a fresh committee until release.  Figure~\ref{fig:shared-coins} gives
the preparation, release, and recovery procedures.
For fixed $U$ and $\boldsymbol{\pk}$, everyone obtains the same ciphertext
vector.  Preparation can take place as soon as the committee has been
assigned.

\begin{figure}[tbp]
\centering
\fbox{%
\begin{minipage}{0.95\textwidth}
\small
\SubprotocolHeading{\Coin}
$\lambda$: the security parameter and coin length.
Use $\ESCR$ from Figure~\ref{fig:conditional-release}, with secret space
$\mathcal M$ and vectors of length $n$.
Fix an encoding $\iota:\bits^\lambda\to\mathcal M$ with efficient
decoding that returns $\bot$ outside its image.
\GadgetHeading{Primitives}
\begin{description}
  \item[$\Coin.\Prepare(U,\boldsymbol{\pk})\to\boldsymbol c$.]
    Input: universal-sampler parameters $U$ and a vector
    $\boldsymbol{\pk}=(\pk_1,\ldots,\pk_n)$ of RNCE public keys.
    Output: a vector $\boldsymbol c=(c_1,\ldots,c_n)$ of ciphertexts.

    Form the description $d_{\boldsymbol{\pk}}$ of the circuit
    \[
      d_{\boldsymbol{\pk}}(q,\rho)
        =\ESCR.\Prepare(\iota(q),\boldsymbol{\pk};\rho).
    \]
    Here $q$ consists of the first $\lambda$ bits of the circuit's random
    input, and $\rho$ supplies the sharing and encryption coins.
    Use a fixed deterministic encoding of this circuit template, with
    $\boldsymbol{\pk}$ hardcoded, and the padding convention of
    Section~\ref{sec:universal-samplers}.
    Return the ciphertext vector
    \[
      \boldsymbol c\gets\Sample^H(U,d_{\boldsymbol{\pk}}).
    \]
    Here $H$ is the universal sampler's public random oracle.

  \item[$\Coin.\Recover(\boldsymbol c,\widehat{\boldsymbol{\sk}})\to Q$.]
    Input: $\boldsymbol c\in(\bits^*)^n$ and
    $\widehat{\boldsymbol{\sk}}\in(\bits^*\cup\{\bot\})^n$.
    Output: an element of $\bits^\lambda\cup\{\bot\}$.
    Compute
    \[
      v\gets\ESCR.\Recover(\boldsymbol c,\widehat{\boldsymbol{\sk}}).
    \]
    If $v=\bot$, return $\bot$; otherwise return the decoding of $v$
    under $\iota$.
\end{description}

\GadgetHeading{Interactions}
\begin{description}
  \item[$\Coin.\Release(\pk,\sk,L_r;\boldsymbol{\pk},\boldsymbol c,B)$.]
    Arguments as for $\ESCR.\Release$.
    Invoke the interaction
    \[
      \ESCR.\Release(\pk,\sk,L_r;
        \boldsymbol{\pk},\boldsymbol c,B).
    \]
\end{description}
\end{minipage}}
\caption{$\Coin$: preparation, release, and recovery of a shared coin.}
\label{fig:shared-coins}
\end{figure}

For coin $Q_s$, the protocol calls $\Coin.\Prepare$ with the reserved
committee's keys and supplies the resulting ciphertext vector to
$\Coin.\Release$.  When using $Q_s$ in a main sampling description,
it fixes all other description fields before enabling release through
$B$.  This test will be specified with the main sampling description.
After the release round, readers form $\widehat{\boldsymbol{\sk}}$
from the ledger as described for $\ESCR$ and obtain $Q_s$ by calling
$\Coin.\Recover$.
